Paste this section into the course template. The tables use booktabs and tabularx. Use the bibliography entries below in the template's bibliography file.

\subsection{Sources and official splits}\label{sources-and-official-splits}

The project used the Exploration-Lab CS-781 Capstone dataset and the
course competition specification \cite{cs781_dataset, cs781_brief}.
Phase 1 supplied Acts and documents from High Courts, the Supreme Court,
and tribunals, together with test prompts. Phase 2 supplied two
annotated legal named entity recognition (NER) folds and a separate test
fold. Phase 3 provided separate training, development, and test splits
for Corporate Law (CL) and Information Technology (IT).

Phases 4 and 5 provided training, development, and test files for court
judgment prediction and legal statute identification. Phase 6 provided
judgment--summary pairs for training and judgments for testing.
Validation judgments were held out from the training pool. Phase 7
supplied separate query files and candidate pools for training,
development, and test. Training and development queries included
relevance labels \cite{cs781_brief}.

\subsection{Corpus statistics}\label{corpus-statistics}

\begin{table}[t]
\centering
\footnotesize
\caption{Phase-wise dataset inventory. The Phase 1 token figure describes the available corpus. The number of tokens used in the pretraining runs is unavailable.}
\label{tab:split-inventory}
\begin{tabularx}{\linewidth}{@{}lXXXX@{}}
\toprule
\textbf{Phase} & \textbf{Task} & \textbf{Training data} & \textbf{Development / validation} & \textbf{Test data} \\
\midrule
1 & Legal language modelling & Four collections, over 4.5 billion tokens available & Run count unavailable & Prompt count unavailable \\
2 & Legal NER & Folds 1 and 2, 12 entity types & No official development split & Fold 3 \\
3 & Rhetorical role prediction & 16,845 sentences & 2,142 sentences & 2,197 sentences \\
4 & Court judgment prediction & 37,387 raw rows, 31,652 unique texts & 3,505 raw rows, 994 clean cases & 1,517 cases \\
5 & Legal statute identification & 42,750 documents & 10,181 documents & 13,019 documents \\
6 & Abstractive summarization & 7,028 usable pairs before validation holdout & Held-out count unavailable & 100 documents \\
7 & Prior case retrieval & Query and candidate counts unavailable & Query and candidate counts unavailable & Query and candidate counts unavailable \\
\bottomrule
\end{tabularx}
\end{table}

\textbf{Phase 2 entity inventory.} The 12 labels are APP, RESP,
A.COUNSEL, R.COUNSEL, JUDGE, WIT, AUTH, COURT, STAT, PREC, DATE, and
CASENO. Fold-level document counts, per-label span counts, and
span-length statistics remain unverified. The training gold annotations
provide the data needed to calculate these distributions
\cite{cs781_brief}.

\textbf{Phase 3 domain and role coverage.} The task used 13 rhetorical
roles. Table~\ref{tab:rr-domains} gives the recorded
document and sentence totals for both domains. Per-role counts within CL
and IT remain unavailable in the run records.

\begin{table}[t]
\centering
\footnotesize
\caption{Phase 3 document and sentence counts by split and domain.}
\label{tab:rr-domains}
\begin{tabularx}{\linewidth}{@{}Xrrrr@{}}
\toprule
\textbf{Split} & \textbf{CL docs} & \textbf{CL sents.} & \textbf{IT docs} & \textbf{IT sents.} \\
\midrule
Training & 40 & 10,593 & 40 & 6,252 \\
Development & 5 & 1,396 & 5 & 746 \\
Test & 5 & 1,339 & 5 & 858 \\
\bottomrule
\end{tabularx}
\end{table}

\textbf{Phase 4 class balance.} The clean development set contained
\textbf{497 class-0 cases and 497 class-1 cases} (50.0\% each). These
counts come from the true-label row sums of the recorded confusion
matrix, \(344+153\) and \(146+351\). The class counts for the training
and private test sets remain unverified.

\textbf{Phase 5 labels and lengths.} The training set contained 100
statute labels and 158,398 label assignments. The median document had
three labels. Individual labels occurred in 272 to 11,731 training
documents.

\begin{table}[t]
\centering
\footnotesize
\caption{Share of Phase 5 documents exceeding each token-length threshold in the recorded audit.}
\label{tab:lsi-lengths}
\begin{tabularx}{\linewidth}{@{}Xrrr@{}}
\toprule
\textbf{Split} & \textbf{>512} & \textbf{>4,096} & \textbf{>8,192} \\
\midrule
Training & 55.3\% & 9.5\% & 2.3\% \\
Development & 60.5\% & 10.8\% & 2.9\% \\
Test & 60.4\% & 9.9\% & 2.5\% \\
\bottomrule
\end{tabularx}
\end{table}

\textbf{Phase 6 document and summary lengths.} The V6 preprocessing code
measures source and reference-summary token lengths before input
selection and target truncation. The completed \texttt{data\_stats.json}
is needed for the actual length distributions, truncation rates, and
holdout size. The 2,048-token context and 768-token target budget are
model settings.

\textbf{Phase 7 retrieval pools.} The preparation code records the
number of queries and candidates per split and the distribution of known
relevant cases per query. The completed \texttt{data\_audit.json} is
needed for those measured values.

\subsection{Preprocessing}\label{preprocessing}

\textbf{Phase 1.} The recorded RoPE and Conv configurations used the
matching custom-trained legal BPE tokenizer,
\texttt{legal\_bpe\_50k\_fixed.json}. The tokenizer audit recorded
\textbf{50,000 vocabulary entries}. The byte-level configuration and the
number of training tokens processed by each model remain unverified.

\textbf{Phase 2.} The data pairs each document ID with its original text
and labelled entity spans. The official prediction format uses
document-level start and end character offsets and an entity label.
Evaluation requires exact agreement on both endpoints and the label. The
implementation used for token-to-character alignment, span decoding, and
entities crossing input windows remains unverified \cite{cs781_brief}.

\textbf{Phase 3.} Each labelled sentence was a classification target.
The context-aware V2 workflow used surrounding sentences and
sentence-position information and preserved the order of predictions
within each document.

\textbf{Phase 4.} The V3.1 run used 28,487 unique texts for fitting and
3,165 for internal validation after auditing repeated records. Long
judgments were represented with up to six 512-token chunks, with
selection weighted toward the end of the document. The model combined
the chunk representations to produce a single judgment prediction.

\textbf{Phase 5.} The reviewed pipeline tokenized complete documents
before chunk selection and represented statute targets as multi-hot
vectors. For overlength training documents, it partitioned the available
chunks into evenly spaced groups and sampled one chunk from each.
Development and test used evenly spaced chunks that included the
beginning and end.

\textbf{Phase 6.} The V6 notebook removed empty reference summaries and
kept duplicate judgments in the same split when creating the holdout. It
reserved up to 768 tokens for the summary within a 2,048-token context.
It used the remaining space for the judgment and prompt markers. The
notebook fitted its passage selector with training references and
prepared a head-and-tail input as an alternative.

\textbf{Phase 7.} The FullFT Pro notebook joined case sentences in
order, preserved IDs as strings, and tokenized queries and candidates
separately. The default represented overlength cases with up to three
384-token chunks spaced across the document. It pooled token
representations into document and local segment vectors. Each training
query could be paired with up to two relevant cases, three hard
negatives from lexical rankings, and one random negative. Retrieval
ranked candidates from the query's own split.
